% Generated mechanically from frozen input p01/ko/conventions.tex.
% Do not edit this generated file; edit the bound locale source instead.

% OK, start here
%

\setcounter{chapter}{1}
\chapter{규약}



\phantomsection
\label{conventions-section-phantom}


\section{설명}
\label{conventions-section-comments}

\noindent
이 문서들을 작성할 때 사용하는 규약의 기본 방침은 실제로 잘 작동하는
규약을 선택하는 것이다.

\section{집합론}
\label{conventions-section-sets}

\noindent
선택 공리를 포함한 Zermelo--Fraenkel 집합론을 사용한다.
\cite{Kunen}을 참조하라. 우주는 사용하지 않는다(이 점은 SGA4와 다르다).
집합론적 문제를 특별히 강조하지는 않지만, 모든 내용이 (물론) 올바르도록
확인하므로 그러한 문제를 무시하지도 않는다.


\section{범주}
\label{conventions-section-categories}

\noindent
범주 $\mathcal{C}$는 대상들의 집합과 각 대상 쌍마다 그 사이에 주어진 사상들의
집합으로 이루어진다. 다시 말해, 다른 문헌에서 ``작은'' 범주라고 부르는
것이다. ``큰'' 범주(대상들이 고유 모임을 이루는 범주)도 사용하지만,
《범주》의 주석 \ref{categories-remark-big-categories}에 열거된 것들만 사용한다.

\section{대수학}
\label{conventions-section-algebra}

\noindent
이 글에서 환은 단위원 $1$을 갖는 가환환을 뜻한다. 따라서 환의 범주는
초기 대상 $\mathbf{Z}$와 끝 대상 $\{0\}$을 갖는다(후자는 $1 = 0$인
유일한 환이다). 가군에서는 단위원이 항등적으로 작용한다고 가정한다.
\cite{Eisenbud}를 참조하라.

\section{표기}
\label{conventions-section-notation}

\noindent
자연수는 $\mathbf{N} = \{1, 2, 3, \ldots\}$의 원소로 정한다.
정수는 $\mathbf{Z} = \{\ldots, -2, -1, 0, 1, 2, \ldots\}$의 원소이다.
유리수체는 $\mathbf{Q}$로 나타낸다.
실수체는 $\mathbf{R}$로 나타낸다.
복소수체는 $\mathbf{C}$로 나타낸다.
